% ============================================================
\chapter{Stochastic Volatility Models}
\label{ch:volatility}
% ============================================================

The Black-Scholes model assumes constant volatility. But markets tell a different story: when one computes implied volatility from observed option prices, it varies with strike and maturity, drawing the celebrated ``volatility smile.'' This phenomenon, which became glaring after the 1987 crash, motivated a new generation of models where volatility itself is a stochastic process. The Heston model (1993), the SABR model (Hagan et al., 2002), and Dupire's local volatility models (1994) each attempt to capture this dynamic.

\section{Motivation: The Volatility Smile}

The Black--Scholes model assumes constant volatility $\sigma$. In practice,
however, the implied volatility $\sigma_{\text{impl}}(K,T)$ depends on the
strike $K$ and maturity $T$, forming a \textbf{volatility surface} with a
characteristic \textit{smile} or \textit{skew}.

\begin{definition}[Volatility Smile]
The \textbf{volatility smile} is the curve
$K \mapsto \sigma_{\text{impl}}(K, T)$ for a fixed maturity.
In equity markets, one typically observes a negative \textit{skew}:
implied volatility is higher for OTM puts (low strikes) than for
OTM calls (high strikes).
\end{definition}

\begin{remark}
The smile is incompatible with the Black--Scholes model and motivates the
introduction of stochastic volatility or jump-diffusion models.
\end{remark}

\section{Local Volatility (Dupire)}

\begin{theorem}[Dupire's Equation (1994)]
If European option prices $C(K,T)$ are known for all $(K,T)$, the
\textbf{local volatility} $\sigma_{\text{loc}}(K,T)$ is defined by:
\[
    \sigma_{\text{loc}}^2(K,T) = \frac{\partial_T C + rK\partial_K C}
    {\frac{1}{2}K^2\partial_{KK}C}.
\]
The underlying then satisfies:
$dS_t = rS_t\,dt + \sigma_{\text{loc}}(S_t, t)S_t\,dW_t^\mathbb{Q}$.
\end{theorem}

\begin{warning}[title=\textbf{Limitations of Local Volatility}]
Dupire's model reproduces today's volatility surface exactly, but its
predictions of future smile dynamics are often unrealistic (the smile
flattens too quickly). Stochastic volatility models offer better dynamics.
\end{warning}

\section{The Heston Model}

\begin{definition}[Heston Model (1993)]
The \textbf{Heston model} describes the underlying and its variance by
the system:
\begin{align}
    dS_t &= rS_t\,dt + \sqrt{v_t}\,S_t\,dW_t^1, \\
    dv_t &= \kappa(\bar{v} - v_t)\,dt + \xi\sqrt{v_t}\,dW_t^2,
\end{align}
with $d\langle W^1, W^2\rangle_t = \rho\,dt$ and parameters:
\begin{itemize}
    \item $v_0 > 0$: initial variance,
    \item $\kappa > 0$: mean-reversion speed,
    \item $\bar{v} > 0$: long-run variance,
    \item $\xi > 0$: volatility of variance (vol of vol),
    \item $\rho \in [-1,1]$: correlation ($\rho < 0$ generates skew).
\end{itemize}
\end{definition}

\begin{theorem}[Feller Condition]
The variance $v_t$ remains strictly positive if and only if the
\textbf{Feller condition} is satisfied:
\[
    2\kappa\bar{v} \geq \xi^2.
\]
\end{theorem}

\subsection{Pricing via Characteristic Function}

\begin{theorem}[Heston Semi-Analytical Formula]
The European call price in the Heston model is:
\[
    C = S_0 P_1 - Ke^{-rT}P_2,
\]
where the probabilities $P_j$ are obtained by Fourier inversion:
\[
    P_j = \frac{1}{2} + \frac{1}{\pi}\int_0^\infty
    \text{Re}\!\left[\frac{e^{-iu\ln K}\varphi_j(u)}{iu}\right]du,
\]
and $\varphi_j(u)$ is the characteristic function of the log-price under
the associated measures.
\end{theorem}

\begin{codeblock}[title=\textbf{Heston Pricing via Fourier Transform}]
\begin{minted}{python}
import numpy as np
from scipy.integrate import quad

def heston_char_func(u, T, r, v0, kappa, vbar, xi, rho):
    """Characteristic function of the log-price under Heston."""
    d = np.sqrt((rho*xi*1j*u - kappa)**2 + xi**2*(1j*u + u**2))
    g = (kappa - rho*xi*1j*u - d) / (kappa - rho*xi*1j*u + d)

    C = r*1j*u*T + (kappa*vbar/xi**2) * (
        (kappa - rho*xi*1j*u - d)*T
        - 2*np.log((1 - g*np.exp(-d*T))/(1 - g))
    )
    D = ((kappa - rho*xi*1j*u - d)/xi**2) * (
        (1 - np.exp(-d*T))/(1 - g*np.exp(-d*T))
    )
    return np.exp(C + D*v0)

def heston_call(S0, K, T, r, v0, kappa, vbar, xi, rho):
    """European call price under the Heston model."""
    ln_K = np.log(K)

    def integrand(u):
        phi = heston_char_func(u - 0.5j, T, r, v0,
                               kappa, vbar, xi, rho)
        return np.real(np.exp(-1j*u*ln_K) * phi / (u**2 + 0.25))

    integral, _ = quad(integrand, 0, 200, limit=200)
    return S0 - np.sqrt(S0*K)*np.exp(-r*T)/np.pi * integral

# Heston parameters
S0, K, T, r = 100, 100, 1.0, 0.05
v0, kappa, vbar, xi, rho = 0.04, 2.0, 0.04, 0.3, -0.7

price = heston_call(S0, K, T, r, v0, kappa, vbar, xi, rho)
print(f"Heston price: {price:.4f}")
\end{minted}
\end{codeblock}

\section{The SABR Model}

\begin{definition}[SABR Model]
The \textbf{SABR} model (Stochastic Alpha Beta Rho, Hagan et al., 2002)
is defined by:
\begin{align}
    dF_t &= \alpha_t F_t^\beta\,dW_t^1, \\
    d\alpha_t &= \nu\alpha_t\,dW_t^2,
\end{align}
with $d\langle W^1, W^2\rangle_t = \rho\,dt$, $F_0 = f$ (forward rate),
$\alpha_0 = \alpha$, and $\beta \in [0,1]$.
\end{definition}

\begin{theorem}[Hagan's Approximation for Implied Volatility]
The Black implied volatility in the SABR model is approximated by:
\[
    \sigma_B(K, f) \approx \frac{\alpha}{(fK)^{(1-\beta)/2}
    \left[1 + \frac{(1-\beta)^2}{24}\ln^2\!\frac{f}{K} + \cdots\right]}
    \cdot \frac{z}{x(z)} \cdot \left[1 + \left(\frac{(1-\beta)^2}{24}
    \frac{\alpha^2}{(fK)^{1-\beta}} + \frac{1}{4}\frac{\rho\beta\nu\alpha}
    {(fK)^{(1-\beta)/2}} + \frac{2-3\rho^2}{24}\nu^2\right)T\right],
\]
where $z = \frac{\nu}{\alpha}(fK)^{(1-\beta)/2}\ln\frac{f}{K}$ and
$x(z) = \ln\!\frac{\sqrt{1 - 2\rho z + z^2} + z - \rho}{1 - \rho}$.
\end{theorem}

\begin{codeblock}[title=\textbf{SABR Implied Volatility}]
\begin{minted}{python}
import numpy as np

def sabr_vol(K, f, T, alpha, beta, rho, nu):
    """Hagan approximation for SABR implied volatility."""
    if abs(f - K) < 1e-10:
        # ATM case
        fk = f**(1 - beta)
        vol = (alpha / fk) * (
            1 + ((1-beta)**2/24 * alpha**2/fk**2
                 + 0.25*rho*beta*nu*alpha/fk
                 + (2-3*rho**2)/24 * nu**2) * T
        )
        return vol

    fk_mid = (f * K)**((1 - beta)/2)
    ln_fk = np.log(f / K)
    z = nu / alpha * fk_mid * ln_fk
    x_z = np.log((np.sqrt(1 - 2*rho*z + z**2) + z - rho)
                 / (1 - rho))

    prefix = alpha / (fk_mid * (1 + (1-beta)**2/24 * ln_fk**2
                                + (1-beta)**4/1920 * ln_fk**4))
    correction = 1 + ((1-beta)**2/24 * alpha**2 / (f*K)**(1-beta)
                      + 0.25*rho*beta*nu*alpha / fk_mid
                      + (2-3*rho**2)/24 * nu**2) * T

    return prefix * z / x_z * correction

# Example: SABR smile
f, T = 0.03, 5.0
alpha, beta, rho, nu = 0.02, 0.5, -0.3, 0.4
strikes = np.linspace(0.01, 0.06, 50)
vols = [sabr_vol(K, f, T, alpha, beta, rho, nu) for K in strikes]
\end{minted}
\end{codeblock}

\section{Calibration}

\begin{definition}[Calibration Problem]
\textbf{Calibration} consists of determining model parameters $\theta^*$
that minimise the discrepancy with market prices (or implied volatilities):
\[
    \theta^* = \arg\min_\theta \sum_{i=1}^{N}
    w_i \bigl(\sigma_{\text{mod}}(K_i, T_i; \theta)
    - \sigma_{\text{mkt}}(K_i, T_i)\bigr)^2.
\]
\end{definition}

\begin{codeblock}[title=\textbf{Heston Model Calibration}]
\begin{minted}{python}
from scipy.optimize import minimize

def calibrate_heston(market_data, S0, r):
    """Calibrate Heston to market implied volatilities.

    market_data: list of (K, T, sigma_mkt)
    """
    def objective(params):
        v0, kappa, vbar, xi, rho = params
        if v0<0 or kappa<0 or vbar<0 or xi<0 or abs(rho)>1:
            return 1e10
        total = 0
        for K, T, sigma_mkt in market_data:
            try:
                C_mod = heston_call(S0, K, T, r, v0, kappa,
                                    vbar, xi, rho)
                sigma_mod = implied_vol_from_price(
                    C_mod, S0, K, T, r)
                total += (sigma_mod - sigma_mkt)**2
            except:
                total += 1e6
        return total

    x0 = [0.04, 1.5, 0.04, 0.3, -0.5]
    bounds = [(0.001,1), (0.01,10), (0.001,1),
              (0.01,2), (-0.99,0.99)]
    result = minimize(objective, x0, method='L-BFGS-B',
                      bounds=bounds)
    return result.x
\end{minted}
\end{codeblock}

\section{Local vs Stochastic Volatility}

\begin{center}
\begin{tabular}{lcc}
    \toprule
    \textbf{Property} & \textbf{Local (Dupire)}
    & \textbf{Stochastic (Heston)} \\
    \midrule
    Exact calibration & Yes & Approximate \\
    Smile dynamics & Unrealistic & Realistic \\
    Number of factors & 1 & 2 \\
    Market & Complete & Incomplete \\
    Greeks & Simple & Complex \\
    \bottomrule
\end{tabular}
\end{center}

\section{Stochastic Local Volatility (SLV) Models}

\begin{definition}[SLV Model]
\textbf{SLV models} combine local and stochastic volatility:
\[
    dS_t = rS_t\,dt + L(S_t, t)\sqrt{v_t}\,S_t\,dW_t^1,
\]
where $L(S,t)$ is the leverage function calibrated to reproduce the
volatility surface exactly, while the dynamics of $v_t$ control the
future behaviour of the smile.
\end{definition}

\section{Exercises}

\begin{exercise}[Volatility Surface]
\begin{enumerate}
    \item Generate option prices in the Heston model for
          $K/S_0 \in \{0.8, 0.9, 1.0, 1.1, 1.2\}$ and $T \in \{0.25, 0.5, 1, 2\}$.
    \item Compute Black--Scholes implied volatility for each price.
    \item Plot the smile for each maturity and comment on the impact of $\rho$.
\end{enumerate}
\end{exercise}

\begin{exercise}[Feller Condition]
\begin{enumerate}
    \item For $\kappa = 2$, $\bar{v} = 0.04$, determine the maximum value
          of $\xi$ satisfying Feller.
    \item Simulate paths of $v_t$ with and without Feller and observe
          behaviour near zero.
\end{enumerate}
\end{exercise}

\begin{exercise}[SABR Calibration]
\begin{enumerate}
    \item Calibrate the SABR model to swaption data (implied volatilities
          for different strikes).
    \item Analyse the stability of parameters $(\alpha, \rho, \nu)$ as a
          function of maturity.
    \item Compare with Heston calibration on the same data.
\end{enumerate}
\end{exercise}

\begin{exercise}[Monte Carlo for Heston]
\begin{enumerate}
    \item Implement the truncated Euler scheme for Heston:
          $v_{t+\Delta t} = \max(v_t + \kappa(\bar{v} - v_t)\Delta t
          + \xi\sqrt{v_t^+}\Delta W^2, 0)$.
    \item Compare with the characteristic function method.
    \item Study the impact of $\rho$ on the smile skew.
\end{enumerate}
\end{exercise}
